% !TEX program = pdflatex
% Hugging Face Deep RL Course - Bonus Unit 1 (complete MDX study edition).
% Compile with latexmk -pdf from notes/en or notes/ru.
\PassOptionsToPackage{T2A}{fontenc}
\PassOptionsToPackage{russian}{babel}
\documentclass[11pt,a4paper]{article}

% Shared typography/layout for Hugging Face Deep RL Course notes.
% Compile with pdfLaTeX / latexmk -pdf with notes/ as the working directory.
% e.g. cd notes && latexmk -pdf unit1.tex

\usepackage[T1]{fontenc}
\usepackage[utf8]{inputenc}
\usepackage[english]{babel}
\usepackage{lmodern}
\usepackage{textcomp}
\usepackage{geometry}
\geometry{a4paper,margin=22mm,headheight=14pt,headsep=8mm,footskip=12mm}

\usepackage{amsmath,amssymb,mathtools}
\usepackage{graphicx,caption,needspace,placeins}
\usepackage{array,booktabs,tabularx}
\usepackage{xcolor}
\usepackage{microtype}
\usepackage{enumitem}
\usepackage{listings}
\usepackage[most]{tcolorbox}
\usepackage{titlesec}
\usepackage{fancyhdr}
\usepackage{hyperref}
\usepackage{bookmark}

\definecolor{ink}{HTML}{243041}
\definecolor{accent}{HTML}{38598C}
\definecolor{secondary}{HTML}{60768A}
\definecolor{codebg}{HTML}{F4F6F8}
\definecolor{codeborder}{HTML}{E3E8EE}
\definecolor{notebg}{HTML}{F3F7FC}

\hypersetup{
  colorlinks=true,
  linkcolor=accent,
  urlcolor=accent,
  citecolor=accent,
  pdfborder={0 0 0}
}
\urlstyle{same}

\setlength{\parindent}{0pt}
\setlength{\parskip}{0.38em plus 0.12em minus 0.05em}

\setlist[itemize]{
  topsep=0.35em,
  itemsep=0.13em,
  leftmargin=1.7em
}
\setlist[enumerate]{
  topsep=0.35em,
  itemsep=0.13em,
  leftmargin=1.8em
}

\setcounter{tocdepth}{2}
\setcounter{secnumdepth}{2}

\titleformat{\section}
  {\Large\bfseries\color{ink}}
  {\thesection}{0.7em}{}

\titleformat{\subsection}
  {\large\bfseries\color{accent}}
  {\thesubsection}{0.7em}{}

\titleformat{\subsubsection}
  {\normalsize\bfseries\color{ink}}
  {\thesubsubsection}{0.6em}{}

\titlespacing*{\section}
  {0pt}{2.2ex plus 0.7ex}{0.9ex}
\titlespacing*{\subsection}
  {0pt}{1.7ex plus 0.4ex}{0.6ex}
\titlespacing*{\subsubsection}
  {0pt}{1.3ex plus 0.3ex}{0.45ex}

\pagestyle{fancy}
\fancyhf{}

\fancyhead[L]{
  \footnotesize\sffamily\color{secondary}
  HF Deep RL Course \enspace/\enspace Unit 1
}
\fancyhead[R]{
  \footnotesize\sffamily\color{secondary}
  \nouppercase{\rightmark}
}
\fancyfoot[C]{
  \footnotesize\color{secondary}\thepage
}
\renewcommand{\headrulewidth}{0.25pt}

\captionsetup{
  font=small,
  labelfont={bf,color=secondary},
  textfont={color=ink},
  justification=centering,
  skip=4pt,
  hypcap=false
}

\lstdefinestyle{hfcode}{
  basicstyle=\ttfamily\footnotesize,
  backgroundcolor=\color{codebg},
  frame=single,
  rulecolor=\color{codeborder},
  numbers=none,
  showstringspaces=false,
  breaklines=true,
  breakatwhitespace=false,
  columns=fullflexible,
  keepspaces=true,
  xleftmargin=0pt,
  xrightmargin=0pt,
  aboveskip=0.6em,
  belowskip=0.6em,
  framexleftmargin=7pt,
  framexrightmargin=7pt,
  framextopmargin=5pt,
  framexbottommargin=5pt,
  upquote=true,
  tabsize=4,
}
\lstset{style=hfcode}

\newcommand{\CourseFigure}[2]{%
  \par\addvspace{0.55em}\noindent%
  \begin{minipage}{\linewidth}%
    \centering
    \includegraphics[
      width=0.86\linewidth,
      height=0.295\textheight,
      keepaspectratio
    ]{#1}\par
    \captionof{figure}{#2}%
  \end{minipage}%
  \par\addvspace{0.55em}%
}

\newtcolorbox{coursenote}{
  enhanced,
  breakable,
  colback=notebg,
  colframe=accent!30,
  boxrule=0.35pt,
  arc=1.2mm,
  left=9pt,
  right=9pt,
  top=6pt,
  bottom=6pt
}

\renewenvironment{quote}
  {\begin{coursenote}}
  {\end{coursenote}}

\emergencystretch=2.5em

\providecommand{\tightlist}{
  \setlength{\itemsep}{0pt}
  \setlength{\parskip}{0pt}
}

\providecommand{\passthrough}[1]{#1}

\Urlmuskip=0mu plus 6mu\relax

\renewcommand{\rmdefault}{cmr}
\renewcommand{\sfdefault}{cmss}
\renewcommand{\ttdefault}{cmtt}
\graphicspath{{../assets/unitbonus1/}}
\fancyhead[L]{\footnotesize\sffamily\color{secondary}Курс Hugging Face Deep RL / Бонусный модуль 1}
\fancyhead[R]{\footnotesize\sffamily\color{secondary}Huggy}
\title{\textbf{Глубокое обучение с подкреплением}\\[0.35em]\large Бонусный модуль 1: Обучаем Huggy и играем с ним}
\author{Курс Hugging Face Deep RL\\\small Полная учебная версия на основе MDX}
\date{}
\begin{document}
\selectlanguage{russian}
\fontencoding{T2A}\selectfont
\maketitle
\vspace{-0.7em}
\begin{center}\small\color{secondary}Источник: курс Hugging Face Deep RL (оригинальный английский текст).\end{center}
\begingroup\small\setlength{\parskip}{0pt}\tableofcontents\endgroup
\vspace{0.5em}
\clearpage

\section{Введение}\label{introduction}

В этом бонусном модуле мы закрепим материал первого модуля: научим собаку Huggy подбирать брошенную палку, а затем сможем \href{https://huggingface.co/spaces/ThomasSimonini/Huggy}{поиграть с обученным агентом прямо в браузере}.

\CourseFigure{thumbnail.png}{Huggy: обложка бонусного модуля курса Hugging Face \protect\href{https://huggingface.co/datasets/huggingface-deep-rl-course/course-images/resolve/main/en/unit2/thumbnail.png}{[оригинал]}}

Итак, приступим!

\section{Как устроен Huggy}\label{how-huggy-works}

Huggy --- среда для глубокого обучения с подкреплением, созданная Hugging Face на основе \href{https://blog.unity.com/technology/puppo-the-corgi-cuteness-overload-with-the-unity-ml-agents-toolkit}{Puppo the Corgi --- проекта команды Unity ML-Agents}. Среда разработана с помощью \href{https://unity.com/}{игрового движка Unity} и \href{https://github.com/Unity-Technologies/ml-agents}{ML-Agents}. ML-Agents --- набор инструментов для Unity, позволяющий \textbf{создавать собственные среды или использовать готовые среды для обучения агентов}.

\CourseFigure{huggy.jpg}{Huggy: собака-агент в среде Unity ML-Agents \protect\href{https://huggingface.co/datasets/huggingface-deep-rl-course/course-images/resolve/main/en/notebooks/unit-bonus1/huggy.jpg}{[оригинал]}}

В этой среде мы хотим научить Huggy \textbf{подбирать брошенную нами палку. Для этого ему необходимо правильно перемещаться в её направлении}.

\subsection{Пространство состояний: что воспринимает Huggy}\label{the-state-space-what-huggy-perceives.}

Huggy не «видит» окружающую среду. Вместо этого мы предоставляем ему следующую информацию:

\begin{itemize}
\tightlist
\item
  Положение цели (палки)
\item
  Относительное положение Huggy и цели
\item
  Ориентация его лап.
\end{itemize}

На основе этой информации Huggy может \textbf{использовать свою политику, чтобы выбрать следующее действие для достижения цели}.

\subsection{Пространство действий: как может двигаться Huggy}\label{the-action-space-what-moves-huggy-can-perform}

\CourseFigure{huggy-action.jpg}{Управление суставными двигателями лап Huggy \protect\href{https://huggingface.co/datasets/huggingface-deep-rl-course/course-images/resolve/main/en/notebooks/unit-bonus1/huggy-action.jpg}{[оригинал]}}

\textbf{Лапы Huggy приводятся в движение суставными двигателями}. Следовательно, чтобы добраться до цели, Huggy должен \textbf{научиться правильно вращать двигатели суставов каждой лапы и таким образом перемещаться}.

\subsection{Функция вознаграждения}\label{the-reward-function}

Функция вознаграждения спроектирована так, чтобы \textbf{Huggy достигал своей цели}: находил палку.

Напомним, что одно из фундаментальных положений обучения с подкреплением --- \emph{гипотеза о вознаграждении}: цель можно описать как \textbf{максимизацию ожидаемого суммарного вознаграждения}.

В нашем случае Huggy должен \textbf{приближаться к палке, не вращаясь слишком сильно}. Поэтому функция вознаграждения должна отражать именно эту цель.

Наша функция вознаграждения:

\CourseFigure{reward.jpg}{Функция вознаграждения для обучения Huggy \protect\href{https://huggingface.co/datasets/huggingface-deep-rl-course/course-images/resolve/main/en/notebooks/unit-bonus1/reward.jpg}{[оригинал]}}

\begin{itemize}
\tightlist
\item
  \emph{Бонус за ориентацию}: мы \textbf{поощряем Huggy за приближение к цели}.
\item
  \emph{Штраф за время}: постоянный штраф на каждом шаге, который \textbf{побуждает его добраться до палки как можно быстрее}.
\item
  \emph{Штраф за вращение}: мы штрафуем Huggy, если \textbf{он слишком сильно вращается и слишком резко поворачивает}.
\item
  \emph{Вознаграждение за достижение цели}: поощряем Huggy за то, что \textbf{он достиг цели}.
\end{itemize}

Если вы хотите увидеть математическое описание этой функции вознаграждения, посмотрите \href{https://blog.unity.com/technology/puppo-the-corgi-cuteness-overload-with-the-unity-ml-agents-toolkit}{презентацию Puppo the Corgi}.

\subsection{Обучение Huggy}\label{train-huggy}

Задача Huggy --- \textbf{научиться правильно и как можно быстрее бежать к цели}. Для этого на каждом шаге, получив наблюдение среды, он должен решить, как вращать суставные двигатели лап, чтобы двигаться правильно (без избыточного вращения) и в направлении цели.

Цикл обучения выглядит следующим образом:

\CourseFigure{huggy-loop.jpg}{Цикл взаимодействия агента Huggy со средой и его обучения \protect\href{https://huggingface.co/datasets/huggingface-deep-rl-course/course-images/resolve/main/en/notebooks/unit-bonus1/huggy-loop.jpg}{[оригинал]}}

Среда обучения выглядит так:

\CourseFigure{training-env.jpg}{Несколько параллельных тренировочных арен со случайным положением палки \protect\href{https://huggingface.co/datasets/huggingface-deep-rl-course/course-images/resolve/main/en/notebooks/unit-bonus1/training-env.jpg}{[оригинал]}}

Здесь \textbf{палка появляется в случайной точке}. Когда Huggy достигает её, палка возникает в другом месте. Для обучения мы создали \textbf{несколько копий среды}. Это ускоряет обучение, поскольку агент получает более разнообразный опыт.

Теперь, когда у вас есть общее представление о среде, вы готовы обучить Huggy подбирать палку.

Для этого мы воспользуемся \href{https://github.com/Unity-Technologies/ml-agents}{ML-Agents}. Не переживайте, если раньше не работали с этим набором инструментов. В данном модуле мы обучим Huggy в Google Colab, а затем вы сможете загрузить своего обученного агента и поиграть с ним прямо в браузере.

В одном из следующих модулей мы подробнее изучим ML-Agents и принципы его работы. А пока упростим задачу и воспользуемся готовой реализацией.

\section{Обучаем Huggy и играем с ним}\label{lets-train-and-play-with-huggy}

\href{https://colab.research.google.com/github/huggingface/deep-rl-class/blob/master/notebooks/bonus-unit1/bonus-unit1.ipynb}{Открыть ноутбук с обучением Huggy в Google Colab}

Мы настоятельно \textbf{рекомендуем студентам выполнять практические задания в Google Colab}, а не на личных компьютерах.

Используя Google Colab, \textbf{можно сосредоточиться на обучении и экспериментах, не отвлекаясь на технические вопросы} настройки окружения.

\subsection{Приступим к обучению Huggy}\label{lets-train-huggy}

\textbf{Чтобы начать обучение Huggy, нажмите кнопку Open in Colab}:

\href{https://colab.research.google.com/github/huggingface/deep-rl-class/blob/master/notebooks/bonus-unit1/bonus-unit1.ipynb}{Открыть в Google Colab}

\CourseFigure{thumbnail.png}{Huggy: обложка бонусного модуля курса Hugging Face \protect\href{https://huggingface.co/datasets/huggingface-deep-rl-course/course-images/resolve/main/en/unit2/thumbnail.png}{[оригинал]}}

В этом ноутбуке мы закрепим материал первого модуля: \textbf{научим Huggy подбирать палку и затем поиграем с обученным агентом прямо в браузере}.

\CourseFigure{huggy.jpg}{Huggy: собака-агент в среде Unity ML-Agents \protect\href{https://huggingface.co/datasets/huggingface-deep-rl-course/course-images/resolve/main/en/notebooks/unit-bonus1/huggy.jpg}{[оригинал]}}

\subsubsection{Среда}\label{the-environment}

\begin{itemize}
\tightlist
\item
  Huggy --- собака-агент в среде, созданной \href{https://twitter.com/ThomasSimonini}{Томасом Симонини} на основе проекта \href{https://blog.unity.com/technology/puppo-the-corgi-cuteness-overload-with-the-unity-ml-agents-toolkit}{Puppo the Corgi}.
\end{itemize}

\subsubsection{Используемая библиотека}\label{the-library-used}

\begin{itemize}
\tightlist
\item
  \href{https://github.com/Unity-Technologies/ml-agents}{MLAgents}
\end{itemize}

Мы постоянно стараемся улучшать наши руководства, поэтому \textbf{если вы обнаружили проблему в этом ноутбуке}, пожалуйста, \href{https://github.com/huggingface/deep-rl-class/issues}{создайте issue в репозитории GitHub}.

\subsection{Цели практического ноутбука}\label{objectives-of-this-notebook}

По окончании работы с ноутбуком вы сможете:

\begin{itemize}
\tightlist
\item
  Разобраться в \textbf{пространстве состояний, пространстве действий и функции вознаграждения}, используемых для обучения Huggy.
\item
  \textbf{Обучить собственного Huggy} подбирать палку.
\item
  \textbf{Играть с обученным Huggy прямо в браузере}.
\end{itemize}

\subsection{Предварительные требования}\label{prerequisites}

Прежде чем приступить к ноутбуку, необходимо:

{[} {]} \textbf{Разобраться с основами обучения с подкреплением} (MC, TD, гипотеза о вознаграждении и т. д.), пройдя первый модуль.

{[} {]} \textbf{Прочитать введение в работу Huggy} в бонусном модуле 1.

\subsection{Настройка GPU}\label{set-the-gpu}

\begin{itemize}
\tightlist
\item
  Для \textbf{ускорения обучения агента мы будем использовать GPU}. Перейдите в \texttt{Runtime\ \textgreater{}\ Change\ Runtime\ type}.
\end{itemize}

\CourseFigure{gpu-step1.jpg}{Google Colab: изменение типа среды выполнения \protect\href{https://huggingface.co/datasets/huggingface-deep-rl-course/course-images/resolve/main/en/notebooks/gpu-step1.jpg}{[оригинал]}}

\begin{itemize}
\tightlist
\item
  Выберите \texttt{Hardware\ Accelerator\ \textgreater{}\ GPU}.
\end{itemize}

\CourseFigure{gpu-step2.jpg}{Google Colab: выбор GPU как аппаратного ускорителя \protect\href{https://huggingface.co/datasets/huggingface-deep-rl-course/course-images/resolve/main/en/notebooks/gpu-step2.jpg}{[оригинал]}}

\subsection{Клонирование репозитория}\label{clone-the-repository}

\begin{itemize}
\tightlist
\item
  Нужно клонировать репозиторий, содержащий \textbf{ML-Agents}.
\end{itemize}

\begin{lstlisting}

# Clone the repository (can take 3min)
git clone --depth 1 https://github.com/Unity-Technologies/ml-agents
\end{lstlisting}

\subsection{Настройка виртуального окружения}\label{setup-the-virtual-environment}

\begin{itemize}
\item
  Для корректной работы \textbf{ML-Agents} в Colab версия Python должна соответствовать требованиям этой библиотеки.
\item
  Поддерживаемую версию Python можно узнать по параметру \texttt{python\_requires} в файлах \texttt{setup.py}. Эти файлы нужны для установки \textbf{ML-Agents}; они расположены по следующим путям:

  \begin{itemize}
  \tightlist
  \item
    \texttt{/content/ml-agents/ml-agents/setup.py}
  \item
    \texttt{/content/ml-agents/ml-agents-envs/setup.py}
  \end{itemize}
\item
  Текущая версия Python в Colab (её можно проверить с помощью \texttt{!python\ -\/-version}) не соответствует параметру \texttt{python\_requires} библиотеки. Из-за этого установка может завершиться без явной ошибки, однако при последующем выполнении команд возникнут такие сообщения:

  \begin{itemize}
  \tightlist
  \item
    \texttt{/bin/bash:\ line\ 1:\ mlagents-learn:\ command\ not\ found}
  \item
    \texttt{/bin/bash:\ line\ 1:\ mlagents-push-to-hf:\ command\ not\ found}
  \end{itemize}
\item
  Для решения проблемы создадим виртуальное окружение с версией Python, совместимой с \textbf{ML-Agents}.
\end{itemize}

\texttt{Примечание:} \emph{Для совместимости с будущими версиями всегда проверяйте параметр \texttt{python\_requires} в установочных файлах. Если версия Python в Colab не поддерживается, укажите в приведённом ниже скрипте максимально поддерживаемую библиотекой версию Python}.

\begin{lstlisting}

# Colab's Current Python Version (Incompatible with ML-Agents)
!python --version
\end{lstlisting}

\begin{lstlisting}

# Install virtualenv and create a virtual environment
!pip install virtualenv
!virtualenv myenv

# Download and install Miniconda
!wget https://repo.anaconda.com/miniconda/Miniconda3-latest-Linux-x86_64.sh
!chmod +x Miniconda3-latest-Linux-x86_64.sh
!./Miniconda3-latest-Linux-x86_64.sh -b -f -p /usr/local

# Activate Miniconda and install Python ver 3.10.12
!source /usr/local/bin/activate
!conda install -q -y --prefix /usr/local python=3.10.12 ujson  # Specify the version here

# Set environment variables for Python and conda paths
!export PYTHONPATH=/usr/local/lib/python3.10/site-packages/
!export CONDA_PREFIX=/usr/local/envs/myenv
\end{lstlisting}

\begin{lstlisting}

# Python Version in New Virtual Environment (Compatible with ML-Agents)
!python --version
\end{lstlisting}

\subsection{Установка зависимостей}\label{installing-the-dependencies}

\begin{lstlisting}

# Go inside the repository and install the package (can take 3min)
%cd ml-agents
pip3 install -e ./ml-agents-envs
pip3 install -e ./ml-agents
\end{lstlisting}

\subsection{Загрузка и распаковка архива среды}\label{download-and-move-the-environment-zip-file-in-.trained-envs-executableslinux}

\begin{itemize}
\tightlist
\item
  Исполняемые файлы нашей среды находятся в ZIP-архиве.
\item
  Нужно скачать архив и поместить его в каталог \texttt{./trained-envs-executables/linux/}.
\end{itemize}

\begin{lstlisting}
mkdir ./trained-envs-executables
mkdir ./trained-envs-executables/linux
\end{lstlisting}

Мы загружаем файл Huggy.zip из \url{https://github.com/huggingface/Huggy} с помощью \texttt{wget}.

\begin{lstlisting}
wget "https://github.com/huggingface/Huggy/raw/main/Huggy.zip" -O ./trained-envs-executables/linux/Huggy.zip
\end{lstlisting}

\begin{lstlisting}
%%capture
unzip -d ./trained-envs-executables/linux/ ./trained-envs-executables/linux/Huggy.zip
\end{lstlisting}

Убедитесь, что файл доступен для запуска.

\begin{lstlisting}
chmod -R 755 ./trained-envs-executables/linux/Huggy
\end{lstlisting}

\subsection{Повторим устройство среды}\label{lets-recap-how-this-environment-works}

\subsubsection{Пространство состояний: что воспринимает Huggy}\label{the-state-space-what-huggy-perceives.-1}

Huggy не «видит» окружающую среду. Вместо этого мы предоставляем ему следующую информацию:

\begin{itemize}
\tightlist
\item
  Положение цели (палки)
\item
  Относительное положение Huggy и цели
\item
  Ориентация его лап.
\end{itemize}

На основе этой информации Huggy \textbf{может выбрать следующее действие для достижения своей цели}.

\CourseFigure{huggy.jpg}{Huggy: собака-агент в среде Unity ML-Agents \protect\href{https://huggingface.co/datasets/huggingface-deep-rl-course/course-images/resolve/main/en/notebooks/unit-bonus1/huggy.jpg}{[оригинал]}}

\subsubsection{Пространство действий: как может двигаться Huggy}\label{the-action-space-what-moves-huggy-can-do}

\CourseFigure{huggy-action.jpg}{Управление суставными двигателями лап Huggy \protect\href{https://huggingface.co/datasets/huggingface-deep-rl-course/course-images/resolve/main/en/notebooks/unit-bonus1/huggy-action.jpg}{[оригинал]}}

\textbf{Лапы Huggy приводятся в движение суставными двигателями}. Чтобы добраться до цели, Huggy должен \textbf{научиться правильно вращать двигатели суставов каждой лапы и таким образом перемещаться}.

\subsubsection{Функция вознаграждения}\label{the-reward-function-1}

Функция вознаграждения спроектирована так, чтобы \textbf{Huggy достигал своей цели}: подбирал палку.

Напомним, что одно из фундаментальных положений обучения с подкреплением --- \emph{гипотеза о вознаграждении}: цель можно описать как \textbf{максимизацию ожидаемого суммарного вознаграждения}.

В нашем случае Huggy должен \textbf{приближаться к палке, не вращаясь слишком сильно}. Поэтому функция вознаграждения должна отражать именно эту цель.

Наша функция вознаграждения:

\CourseFigure{reward.jpg}{Функция вознаграждения для обучения Huggy \protect\href{https://huggingface.co/datasets/huggingface-deep-rl-course/course-images/resolve/main/en/notebooks/unit-bonus1/reward.jpg}{[оригинал]}}

\begin{itemize}
\tightlist
\item
  \emph{Бонус за ориентацию}: мы \textbf{поощряем Huggy за приближение к цели}.
\item
  \emph{Штраф за время}: постоянный штраф на каждом шаге, который \textbf{побуждает его добраться до палки как можно быстрее}.
\item
  \emph{Штраф за вращение}: мы штрафуем Huggy, если \textbf{он слишком сильно вращается и слишком резко поворачивает}.
\item
  \emph{Вознаграждение за достижение цели}: поощряем Huggy за то, что \textbf{он достиг цели}.
\end{itemize}

\subsection{Проверка конфигурации Huggy}\label{check-the-huggy-config-file}

\begin{itemize}
\item
  В ML-Agents \textbf{гиперпараметры обучения задаются в файлах config.yaml}.
\item
  В рамках этого ноутбука мы не будем изменять гиперпараметры. Но если вы захотите поэкспериментировать, Unity предоставляет \href{https://github.com/Unity-Technologies/ml-agents/blob/main/docs/Training-Configuration-File.md}{подробную документацию по каждому из них}.
\item
  Нужно создать конфигурационный файл для Huggy.
\item
  Перейдите в каталог \texttt{/content/ml-agents/config/ppo}.
\item
  Создайте новый файл \texttt{Huggy.yaml}.
\item
  Скопируйте и вставьте приведённое ниже содержимое.
\end{itemize}

\begin{lstlisting}
behaviors:
  Huggy:
    trainer_type: ppo
    hyperparameters:
      batch_size: 2048
      buffer_size: 20480
      learning_rate: 0.0003
      beta: 0.005
      epsilon: 0.2
      lambd: 0.95
      num_epoch: 3
      learning_rate_schedule: linear
    network_settings:
      normalize: true
      hidden_units: 512
      num_layers: 3
      vis_encode_type: simple
    reward_signals:
      extrinsic:
        gamma: 0.995
        strength: 1.0
    checkpoint_interval: 200000
    keep_checkpoints: 15
    max_steps: 2e6
    time_horizon: 1000
    summary_freq: 50000
\end{lstlisting}

\begin{itemize}
\item
  Не забудьте сохранить файл!
\item
  \textbf{Если вы захотите изменить гиперпараметры}, в ноутбуке Google Colab можно открыть файл config.yaml по пути \texttt{/content/ml-agents/config/ppo/Huggy.yaml}.
\end{itemize}

Теперь всё готово к обучению агента.

\subsection{Обучение агента}\label{train-our-agent}

Чтобы обучить агента, достаточно \textbf{запустить mlagents-learn и указать исполняемый файл среды}.

\CourseFigure{mllearn.png}{Структура команды mlagents-learn и её параметры \protect\href{https://huggingface.co/datasets/huggingface-deep-rl-course/course-images/resolve/main/en/notebooks/unit-bonus1/mllearn.png}{[оригинал]}}

В ML-Agents обучение запускается отдельным скриптом. Мы задаём четыре параметра:

\begin{enumerate}
\def\labelenumi{\arabic{enumi}.}
\tightlist
\item
  \texttt{mlagents-learn\ \textless{}config\textgreater{}}: путь к конфигурационному файлу с гиперпараметрами.
\item
  \texttt{-\/-env}: расположение исполняемого файла среды.
\item
  \texttt{-\/-run-id}: имя запуска обучения.
\item
  \texttt{-\/-no-graphics}: отключение графического отображения во время обучения.
\end{enumerate}

Обучите модель. Если выполнение прервётся, используйте флаг \texttt{-\/-resume} для продолжения обучения.

\begin{quote}
При первом использовании \texttt{-\/-resume} возникнет ошибка; повторно выполните блок, чтобы её обойти.
\end{quote}

Обучение займёт от 30 до 45 минут в зависимости от компьютера (не забудьте \textbf{настроить GPU}). Можно пока выпить заслуженную чашку кофе.

\begin{lstlisting}
mlagents-learn ./config/ppo/Huggy.yaml --env=./trained-envs-executables/linux/Huggy/Huggy --run-id="Huggy" --no-graphics
\end{lstlisting}

\subsection{Публикация агента на Hugging Face Hub}\label{push-the-agent-to-the-hub}

\begin{itemize}
\tightlist
\item
Мы обучили агента и теперь \textbf{готовы загрузить его на Hub, чтобы играть с Huggy в браузере}.
\end{itemize}

Чтобы поделиться своей моделью с сообществом, осталось выполнить ещё три шага:

\begin{enumerate}
\def\labelenumi{\arabic{enumi}.}
\item
  Если аккаунта ещё нет, зарегистрируйтесь на Hugging Face: \url{https://huggingface.co/join}.
\item
  Войдите в аккаунт и получите токен доступа на сайте Hugging Face.
\end{enumerate}

\begin{itemize}
\tightlist
\item
  Создайте новый токен (\url{https://huggingface.co/settings/tokens}) \textbf{с правами записи (write)}.
\end{itemize}

\CourseFigure{create-token.jpg}{Создание токена Hugging Face с правом записи \protect\href{https://huggingface.co/datasets/huggingface-deep-rl-course/course-images/resolve/main/en/notebooks/create-token.jpg}{[оригинал]}}

\begin{itemize}
\tightlist
\item
  Скопируйте токен.
\item
  Запустите следующую ячейку и вставьте токен.
\end{itemize}

\begin{lstlisting}
from huggingface_hub import notebook_login

notebook_login()
\end{lstlisting}

Если вы не используете Google Colab или Jupyter Notebook, вместо этого выполните команду \texttt{huggingface-cli\ login}.

После этого достаточно запустить \texttt{mlagents-push-to-hf}.

\CourseFigure{mlpush.png}{Структура команды mlagents-push-to-hf \protect\href{https://huggingface.co/datasets/huggingface-deep-rl-course/course-images/resolve/main/en/notebooks/unit-bonus1/mlpush.png}{[оригинал]}}

Для этой команды указываются четыре параметра:

\begin{enumerate}
\def\labelenumi{\arabic{enumi}.}
\tightlist
\item
  \texttt{-\/-run-id}: идентификатор запуска обучения.
\item
  \texttt{-\/-local-dir}: каталог, в котором сохранён агент, то есть \texttt{results/\textless{}run\_id\ name\textgreater{}}; в моём случае --- results/First Training.
\item
  \texttt{-\/-repo-id}: имя репозитория Hugging Face, который нужно создать или обновить. Оно имеет формат \texttt{\textless{}your\ huggingface\ username\textgreater{}/\textless{}the\ repo\ name\textgreater{}}. Если репозитория ещё нет, \textbf{он будет создан автоматически}.
\item
  \texttt{-\/-commit-message}: сообщение коммита, поскольку репозитории Hugging Face используют Git.
\end{enumerate}

\begin{lstlisting}
mlagents-push-to-hf --run-id="HuggyTraining" --local-dir="./results/Huggy" --repo-id="ThomasSimonini/ppo-Huggy" --commit-message="Huggy"
\end{lstlisting}

Если всё прошло успешно, в конце процесса появится примерно такое сообщение (но с другим URL):

\begin{lstlisting}
Your model is pushed to the hub. You can view your model here: https://huggingface.co/ThomasSimonini/ppo-Huggy
\end{lstlisting}

Это ссылка на репозиторий вашей модели. В нём находятся карточка модели с инструкцией по использованию, логи TensorBoard и конфигурационный файл. \textbf{Поскольку это Git-репозиторий, вы можете создавать разные коммиты, обновлять модель повторными публикациями, открывать Pull Request и т. д.}

\CourseFigure{modelcard.png}{Пример репозитория обученной модели Huggy и карточки Model Card \protect\href{https://huggingface.co/datasets/huggingface-deep-rl-course/course-images/resolve/main/en/notebooks/unit-bonus1/modelcard.png}{[оригинал]}}

А теперь самое интересное: \textbf{можно играть с Huggy онлайн}.

\subsection{Играем со своим Huggy}\label{play-with-your-huggy}

Этот шаг самый простой:

\begin{itemize}
\item
  Откройте игру Huggy в браузере: \url{https://huggingface.co/spaces/ThomasSimonini/Huggy}.
\item
  Нажмите кнопку Play with my Huggy model.
\end{itemize}

\CourseFigure{load-huggy.jpg}{Браузерная демонстрация Huggy: выбор модели из аккаунта \protect\href{https://huggingface.co/datasets/huggingface-deep-rl-course/course-images/resolve/main/en/notebooks/unit-bonus1/load-huggy.jpg}{[оригинал]}}

\begin{enumerate}
\def\labelenumi{\arabic{enumi}.}
\item
  На шаге 1 введите своё имя пользователя (регистр символов важен: например, мой ник --- ThomasSimonini, а не thomassimonini или ThOmasImoNInI) и нажмите кнопку поиска.
\item
  На шаге 2 выберите репозиторий своей модели.
\item
  На шаге 3 \textbf{выберите модель для воспроизведения}:
\end{enumerate}

\begin{itemize}
\tightlist
\item
  У меня их несколько, поскольку мы сохраняли модель каждые 500000 шагов.
\item
  Но для получения самого свежего результата я выбираю \texttt{Huggy.onnx}.
\end{itemize}

Полезно \textbf{попробовать модели с разным числом шагов обучения, чтобы увидеть прогресс агента}.

Поздравляем с завершением этого бонусного модуля!

Теперь можно отдохнуть и поиграть со своим Huggy. И \textbf{не забудьте поделиться Huggy с друзьями}! Если будете рассказывать о нём в социальных сетях, \textbf{отметьте @huggingface и @simoninithomas}.

\CourseFigure{huggy-cover.jpeg}{Huggy в игровой среде \protect\href{https://huggingface.co/datasets/huggingface-deep-rl-course/course-images/resolve/main/en/notebooks/unit-bonus1/huggy-cover.jpeg}{[оригинал]}}

\subsection{Продолжайте учиться и открывать новое}\label{keep-learning-stay-awesome}

\section{Игра с Huggy}\label{play-with-huggy}

Теперь, когда вы обучили Huggy и опубликовали его на Hub, \textbf{можно поиграть с ним}.

Здесь всё просто:

\begin{itemize}
\tightlist
\item
  Откройте игру Huggy в браузере: \url{https://huggingface.co/spaces/ThomasSimonini/Huggy}.
\item
  Нажмите кнопку Play with my Huggy model.
\end{itemize}

\CourseFigure{load-huggy.jpg}{Браузерная демонстрация Huggy: выбор модели из аккаунта \protect\href{https://huggingface.co/datasets/huggingface-deep-rl-course/course-images/resolve/main/en/notebooks/unit-bonus1/load-huggy.jpg}{[оригинал]}}

\begin{enumerate}
\def\labelenumi{\arabic{enumi}.}
\item
  На шаге 1 выберите репозиторий модели --- её идентификатор (в моём случае ThomasSimonini/ppo-Huggy).
\item
  На шаге 2 \textbf{выберите модель для воспроизведения}:
\end{enumerate}

\begin{itemize}
\tightlist
\item
  У меня их несколько, поскольку мы сохраняли модель каждые 500000 шагов.
\item
  Если нужна самая свежая версия, я выбираю Huggy.onnx.
\end{itemize}

Полезно \textbf{пробовать разные контрольные точки модели, чтобы наблюдать, как улучшается агент}.

\section{Заключение}\label{conclusion}

Поздравляем с завершением этого бонусного модуля!

Теперь можно отдохнуть и поиграть со своим Huggy. И \textbf{не забудьте поделиться Huggy с друзьями}! Если будете рассказывать о нём в социальных сетях, \textbf{отметьте @huggingface и @simoninithomas}.

\CourseFigure{huggy-cover.jpeg}{Huggy в игровой среде \protect\href{https://huggingface.co/datasets/huggingface-deep-rl-course/course-images/resolve/main/en/notebooks/unit-bonus1/huggy-cover.jpeg}{[оригинал]}}

Наконец, нам очень важно \textbf{узнать ваше мнение о курсе и о том, как его улучшить}. Если у вас есть пожелания, пожалуйста, \href{https://forms.gle/BzKXWzLAGZESGNaE9}{заполните эту форму}.

\subsubsection{Продолжайте учиться и открывать новое}\label{keep-learning-stay-awesome-1}

\end{document}
